\documentclass[11pt]{article}
\usepackage[margin=1.1in]{geometry}
\usepackage{amsmath,amssymb,amsthm}
\usepackage{natbib}
\usepackage{booktabs}
\usepackage{hyperref}

\newtheorem{proposition}{Proposition}
\newtheorem{lemma}[proposition]{Lemma}
\theoremstyle{definition}
\newtheorem{definition}{Definition}

\newcommand{\R}{\mathbb{R}}
\newcommand{\E}{\mathbb{E}}
\DeclareMathOperator*{\argmax}{arg\,max}
\newcommand{\key}[1]{\texttt{#1}}

\title{Editing in Print}
\author{Rikuy\thanks{A typeset editor for \LaTeX{} sources. Everything on this page is live.}}
\date{\today}

\begin{document}
\maketitle

\begin{abstract}
You are reading a \LaTeX{} document that you can edit where it stands. Click into any sentence and type. Each keystroke becomes a minimal patch to the source, so the \texttt{.tex} file stays yours: no import, no export, no lock in.
\end{abstract}

\section{How to use it}
Click anywhere in this paragraph and start typing. Select a phrase and press \key{Ctrl+B} to make it \textbf{bold} or \key{Ctrl+I} to \emph{emphasise} it. Press \key{Enter} for a new paragraph. Typing a percent sign gives you 50\%, not a comment, because Rikuy escapes special characters for you.

Math is an object on the page. Click $\alpha + \beta = \gamma$ to edit its source with a live preview, or type a dollar sign to insert new math. Display equations work the same way, numbers and all:
\begin{equation}
  V(k) = \max_{c \in [0,k]} \bigl\{ u(c) + \beta\, V\bigl(f(k) - c\bigr) \bigr\}. \label{eq:bellman}
\end{equation}
References follow their targets: equation~\eqref{eq:bellman} renumbers itself if you add an equation above it.

\section{The round trip}\label{sec:roundtrip}
The page and the source never disagree, because the page is computed from the source and every edit is computed back into it.

\begin{definition}[Source map]
A \emph{source map} assigns to every visible character on the page the exact span of source text that produced it.
\end{definition}

\begin{proposition}[Minimal patches]\label{prop:patch}
If an edit only touches characters inside one source map span, the patch changes only that span. Markup outside the selection is never rewritten.
\end{proposition}

\begin{proof}
Rikuy never regenerates the file. Insertions are escaped and placed at the mapped offset; deletions that would cut through markup fall back to removing only the visible characters, which leaves every brace balanced.
\end{proof}

Proposition~\ref{prop:patch} is why your diffs stay small:\footnote{Try it: edit this footnote, then look at the source pane.} the history of a paper edited here reads like one edited by hand.

\section{Citations and cross references}
Citations come from the bibliography file in this project. \citet{knuth1984} designed \TeX{} around a source language, and \citet{lamport1994} built \LaTeX{} on top of it. Editing the typeset output directly is an old idea \citep{kastrup2002, wen2018}, usually defeated by the gap between source and print. Section~\ref{sec:roundtrip} explains how Rikuy closes it.
\begin{align}
  \E[\,y \mid x\,] &= \int_{\R} y \, f(y \mid x)\, dy \label{eq:cond}\\
  \hat\theta &= \argmax_{\theta \in \Theta} \sum_{i=1}^{n} \log f(y_i \mid x_i; \theta) \label{eq:mle}
\end{align}
Both rows of the alignment above carry their own numbers, (\ref{eq:cond}) and (\ref{eq:mle}).

\section{Keyboard}
Table~\ref{tab:keys} lists the shortcuts. Every cell is editable.

\begin{table}[h]
\centering
\begin{tabular}{ll}
\toprule
Keys & Action \\
\midrule
Ctrl+B, Ctrl+I & Bold, emphasis \\
Ctrl+M or \$ & New inline math \\
\textbackslash & Insert a command \\
Ctrl+Z, Ctrl+Shift+Z & Undo, redo \\
Alt+click & Show in source \\
\bottomrule
\end{tabular}
\caption{Shortcuts. On a Mac, use Cmd for Ctrl.}\label{tab:keys}
\end{table}

\section{The exact PDF}
The pages you edit are typeset by the browser, so their line and page breaks are close to \TeX{} but not identical: the browser breaks lines one at a time, while \TeX{} chooses the breaks for a whole paragraph at once \citep{knuthplass1981}. The PDF button runs the real thing: pdf\LaTeX{} from \TeX{} Live 2026, compiled to WebAssembly and running in this browser. Its PDF is the one \texttt{pdflatex} would give you, it recompiles a moment after each edit, and a click on it finds the same sentence here. Packages that draw, such as TikZ, show as source blocks on these pages and fully drawn in the PDF.

\begin{itemize}
  \item Open a \texttt{.tex}, a \texttt{.bib}, or an Overleaf \texttt{.zip}.
  \item Save downloads the whole project as a \texttt{.zip}.
  \item Your work is kept in this browser between visits.
\end{itemize}

\bibliographystyle{plainnat}
\bibliography{refs}
\end{document}
